% tab_features.tex: generated by paper2-tnnls/scripts/tab_features.py from committed runs/ files. Do not edit by hand.
% Study 2 feature groups by formula; no feature is described by any role it has in another analysis. Use table* or a smaller font; the Definition column wraps.
\begin{tabular}{@{}lcp{0.30\linewidth}p{0.42\linewidth}@{}}
\toprule
Group & Count & Features & Definition \\
\midrule
% run: analysis_posthoc/learned_ceiling/data.py feature_names(), engineer() (names and formulas); runs/P2-LRF-FROZEN-LFIX/spec.json features (order, n_features = 35); group counts computed
token scores & 7 & \texttt{B}, \texttt{D}, \texttt{q}, \texttt{s3}, \texttt{tbv}, \texttt{ltex}, \texttt{nbr} & $B_j$: log-probability of the observed code at token $j$ with the tokens of $j$'s group masked, the other image tokens visible and the caption masked; $D_j$: the same with every token masked; $q_j=B_j-D_j$; \texttt{s3} $=1-\max_y p(y)$ of the pass that gives $B_j$; \texttt{tbv}: log-probability of the observed code from one raw-logit pass with all image tokens visible and the caption masked; \texttt{ltex} $=\log(1+\tau_j)$, $\tau_j$ the pixel variance of the $16\times16$ patch of token $j$; \texttt{nbr}: mean $L_2$ distance of the codebook vector of token $j$ to those of its eight neighbours \\
% run: analysis_posthoc/learned_ceiling/data.py feature_names(), engineer() (names and formulas); runs/P2-LRF-FROZEN-LFIX/spec.json features (order, n_features = 35); group counts computed
position & 3 & \texttt{row}, \texttt{col}, \texttt{rad} & \texttt{row}, \texttt{col}: grid row and column of the token divided by 31; \texttt{rad} $=\sqrt{(\mathrm{row}-\tfrac12)^2+(\mathrm{col}-\tfrac12)^2}$ \\
% run: analysis_posthoc/learned_ceiling/data.py feature_names(), engineer() (names and formulas); runs/P2-LRF-FROZEN-LFIX/spec.json features (order, n_features = 35); group counts computed
neighbourhood & 8 & \texttt{nb3mean\_q}, \texttt{nb3mean\_B}, \texttt{nb3mean\_s3}, \texttt{nb3mean\_tbv}, \texttt{nb3mean\_ltex}, \texttt{nb3mean\_nbr}, \texttt{nb3min\_q}, \texttt{nb5mean\_q} & \texttt{nb3mean\_}$k$: mean of score $k$ over the eight neighbours (edge-aware, token excluded); \texttt{nb3min\_q}: their minimum of $q$; \texttt{nb5mean\_q}: mean of $q$ over the $5\times5$ neighbourhood without the token \\
% run: analysis_posthoc/learned_ceiling/data.py feature_names(), engineer() (names and formulas); runs/P2-LRF-FROZEN-LFIX/spec.json features (order, n_features = 35); group counts computed
within-image rank & 4 & \texttt{rank\_q}, \texttt{rank\_B}, \texttt{rank\_s3}, \texttt{rank\_tbv} & percentile rank of the score among the 1\,024 tokens of the image, in $[0,1]$ \\
% run: analysis_posthoc/learned_ceiling/data.py feature_names(), engineer() (names and formulas); runs/P2-LRF-FROZEN-LFIX/spec.json features (order, n_features = 35); group counts computed
image mean & 6 & \texttt{itemmean\_q}, \texttt{itemmean\_B}, \texttt{itemmean\_s3}, \texttt{itemmean\_tbv}, \texttt{itemmean\_ltex}, \texttt{itemmean\_nbr} & mean of the score over the 1\,024 tokens of the image \\
% run: analysis_posthoc/learned_ceiling/data.py feature_names(), engineer() (names and formulas); runs/P2-LRF-FROZEN-LFIX/spec.json features (order, n_features = 35); group counts computed
image standard deviation & 6 & \texttt{itemsd\_q}, \texttt{itemsd\_B}, \texttt{itemsd\_s3}, \texttt{itemsd\_tbv}, \texttt{itemsd\_ltex}, \texttt{itemsd\_nbr} & standard deviation of the score over the 1\,024 tokens of the image \\
% run: analysis_posthoc/learned_ceiling/data.py feature_names(), engineer() (names and formulas); runs/P2-LRF-FROZEN-LFIX/spec.json features (order, n_features = 35); group counts computed
image z-score & 1 & \texttt{itemz\_q} & $(q_j-\bar q)/(\mathrm{sd}(q)+10^{-6})$ within the image \\
% run: sum of the group counts (asserted equal to 35)
\textbf{Total} & 35 &  &  \\
\bottomrule
\end{tabular}
